\begin{helper}
Let \(s,k,a,b\) be natural numbers.  If
\(\mathrm{RamseyProperty}(s,k+1,a)\) and
\(\mathrm{RamseyProperty}(s+1,k,b)\) both hold, then
\(\mathrm{RamseyProperty}(s+1,k+1,a+b+1)\) holds.
\end{helper}
\begin{proof}
Let \(G\) be an arbitrary simple graph on a vertex set of size \(a+b+1\),
identified with \(\operatorname{Fin}(a+b+1)\).  Fix one vertex \(v\).  Among
the other \(a+b\) vertices, let \(N\) be the set of neighbors of \(v\) and let
\(M\) be the set of non-neighbors of \(v\).  These two sets partition the
vertices other than \(v\), so \(|N|+|M|=a+b\).

First suppose \(|N|\ge a\).  Choose a subset \(A\subseteq N\) with
\(|A|=a\), and identify \(A\) with \(\operatorname{Fin}(a)\) by a bijection.
Pulling \(G\) back along this bijection gives a graph on
\(\operatorname{Fin}(a)\).  By the hypothesis
\(\mathrm{RamseyProperty}(s,k+1,a)\), as interpreted by
\(\noderef{RamseyProperty}\), this restricted graph has either an \(s\)-clique
or a \((k+1)\)-clique in its complement.  In the first case, map the
\(s\)-clique back into \(A\).  Since every vertex of \(A\) is adjacent to
\(v\), adjoining \(v\) produces an \((s+1)\)-clique in \(G\).  In the second
case, mapping the complement clique back into \(A\) preserves non-adjacency
and cardinality, so it is already a \((k+1)\)-clique in \(G^{c}\).

It remains to consider the case \(|N|<a\).  From \(|N|+|M|=a+b\) it follows
that \(|M|\ge b\).  Choose a subset \(B\subseteq M\) with \(|B|=b\), identify
\(B\) with \(\operatorname{Fin}(b)\), and pull \(G\) back to a graph on
\(\operatorname{Fin}(b)\).  By the hypothesis
\(\mathrm{RamseyProperty}(s+1,k,b)\), there is either an \((s+1)\)-clique in
this restricted graph or a \(k\)-clique in its complement.  In the first case,
mapping the clique back into \(B\) gives an \((s+1)\)-clique in \(G\).  In the
second case, mapping the \(k\)-clique back into \(B\) gives a \(k\)-clique in
\(G^{c}\); because each vertex of \(B\) is a non-neighbor of \(v\), adjoining
\(v\) gives a \((k+1)\)-clique in \(G^{c}\).

Thus in all cases \(G\) contains either an \((s+1)\)-clique or a
\((k+1)\)-clique in its complement.  Since \(G\) was arbitrary, this is exactly
\(\mathrm{RamseyProperty}(s+1,k+1,a+b+1)\).
\end{proof}
